\section{Results}
\label{sec:results}

\subsection{Proxy Similarity Is Necessary but Not Sufficient}
\label{sec:proxy_results}

The first result is scoped to proxy preservation. Across four single-object Isaac Sim
assets, converted and source MDL views stay close by proxy metrics. Mean PSNR
is 35.52\,dB, SSIM 0.990, LPIPS 0.020, CLIP 0.925, and DINOv2 0.872 over 24
matched pairs. Those numbers pass this visual check while leaving grounding
behavior unresolved. A close render can still alter the target, point frame, or
material cue. We use proxy match as a prerequisite for later task checks. The
render panel gives context; reported numbers come from frozen
tables.

\begin{figure*}[t]
\centering
\includegraphics[width=\textwidth]{figures/fig_render_scene_evidence_wide.png}
\caption{Actual render examples for the measured tasks. Top: a representative
cropped single-object proxy detail pair. Bottom: a representative GRScenes
original/noMDL commercial-interior render pair. These panels orient readers;
metric and task measurements come from the frozen tables.}
\label{fig:render_scene_evidence_acl}
\end{figure*}

\subsection{VLM Grounding Under a Frozen Material Perturbation}
\label{sec:vlm_grounding_pilot}

The clean GRScenes pool has 15 source/noMDL pairs that pass QA. It serves as a
pilot preservation control because it remains below the configured 20-pair clean
final set. The lower row of
Figure~\ref{fig:render_scene_evidence_acl} anchors the section in actual
commercial-interior renders; reported numbers remain tied to frozen pair tables.

\begin{table*}[t]
\centering
\caption{Clean-pool GRScenes VLM grounding pilot over 15 blind-visual-QA PASS original/converted pairs. Counts are O/C hits under the structured-text prompt; the pool remains below the 20-pair benchmark threshold and is a control, not final benchmark evidence.}
\label{tab:grscenes_vlm_clean_pool_pass15}
\footnotesize
\setlength{\tabcolsep}{4pt}
\begin{tabular}{@{}>{\raggedright\arraybackslash}p{0.15\textwidth}
                >{\raggedright\arraybackslash}p{0.07\textwidth}
                >{\raggedright\arraybackslash}p{0.12\textwidth}
                >{\raggedright\arraybackslash}p{0.12\textwidth}
                >{\raggedright\arraybackslash}p{0.13\textwidth}
                >{\raggedright\arraybackslash}p{0.20\textwidth}@{}}
\toprule
\textbf{Probe} & \textbf{Rows} & \textbf{Answer} & \textbf{Raw point} & \textbf{Norm-1000 point} & \textbf{Scope} \\
\midrule
Gemma4 local & 30/30 & 15/15 / 15/15 & 0/15 / 0/15 & 8/15 / 6/15 & pilot preservation control \\
Qwen2.5-VL structured & 23/30 & 8/11 / 9/12 & 5/14 / 5/15 & 0/14 / 0/15 & coordinate-format diagnostic \\
\bottomrule
\end{tabular}
\par\vspace{0.25em}
\begin{minipage}{0.88\textwidth}
\scriptsize
\emph{Notes.} Gemma4 has 11/15 normalized-1000 same-hit-status pairs, 5/15
both-hit pairs, and 27.90 px mean point delta. Qwen2.5-VL normalized agreement
is an all-miss agreement over comparable points; raw pair agreement is 9/14
with 2/14 both-hit.
\end{minipage}
\end{table*}


Table~\ref{tab:grscenes_vlm_clean_pool_pass15} separates answer stability from
point-grounding stability. Gemma4 answers all 30 clean-pool renders correctly,
with normalized-1000 point hits of 8/15 original and 6/15 converted. Qwen2.5-VL
has lower parser coverage: 23/30 scorable answer rows, raw point hits of
5/14 and 5/15, and zero normalized-1000 point hits in both conditions.

The expanded stress set moves to the task question: when the target stays
visible but the material channel can move, does the same prompt-and-coordinate format still
bind model output to the same object? We freeze 30 target-centered PASS/WARN
pairs for this test.

\begin{table*}[t]
\centering
\caption{Expanded GRScenes material-shift grounding results over a frozen 30-pair target-centered stress set. Gemma4 is the canonical root run for this frozen set; Qwen2.5-VL is a second-model diagnostic under the same manifest. Values are hit counts for original/converted renders unless otherwise noted. Norm-1000 point-in-bbox is the primary point metric; raw point-in-bbox is retained only as a coordinate diagnostic.}
\label{tab:grscenes_vlm_stress_expanded30}
\scriptsize
\resizebox{\textwidth}{!}{%
\begin{tabular}{lllccccc}
\toprule
\textbf{Model} & \textbf{Role} & \textbf{Answer rows} & \textbf{Answer O/C} & \textbf{Raw point O/C} & \textbf{Norm-1000 point O/C} & \textbf{Norm hit agree} & \textbf{Norm both-hit} \\
\midrule
Gemma4 local & canonical & 60/60 & 30/30 / 30/30 & 1/30 / 1/30 & 27/30 / 29/30 & 28/30 & 27/30 \\
Qwen2.5-VL & diagnostic & 55/60 & 27/29 / 24/26 & 22/29 / 19/29 & 3/29 / 3/29 & 24/28 & 1/28 \\
\bottomrule
\end{tabular}%
}
\end{table*}


Table~\ref{tab:grscenes_vlm_stress_expanded30} gives the main result for this
frozen stress test. Gemma4 keeps all 30 target labels in both material
conditions. Under the normalized-1000 scoring format, it scores 27/30 point hits on
original renders and 29/30 on converted renders; 28 pairs keep the same hit
status. Qwen2.5-VL has another coordinate preference: raw-pixel hits are 22/29
original and 19/29 converted, while normalized-1000 hits stay 3/29 for both
versions. The comparison should be read as a protocol diagnostic: material
shift, prompt schema, and coordinate frame meet at the measurement interface.

\begin{table*}[t]
\centering
\caption{Reviewer-closure paired bootstrap confidence intervals for existing VLM and InternNav evidence. Deltas are converted/noMDL minus original MDL and are descriptive within the current frozen evidence pools, not population-level guarantees.}
\label{tab:reviewer_closure_paired_ci}
\scriptsize
\resizebox{\textwidth}{!}{%
\begin{tabular}{lllrrrr}
\toprule
\textbf{Evidence} & \textbf{Metric} & \textbf{N} & \textbf{Original} & \textbf{Converted} & \textbf{Delta [95\% CI]} & \textbf{Direction} \\
\midrule
GRScenes VLM expanded30 & Gemma4 Norm-1000 point hit & 30 & 0.9 & 0.9667 & 0.0667 [0.0, 0.1667] & converted higher \\
GRScenes VLM expanded30 & Gemma4 Raw point diagnostic & 30 & 0.0333 & 0.0333 & 0.0 [0.0, 0.0] & tie \\
GRScenes VLM expanded30 & Qwen2.5-VL Norm-1000 point hit & 28 & 0.1071 & 0.1071 & 0.0 [-0.1429, 0.1429] & tie \\
GRScenes VLM expanded30 & Qwen2.5-VL Raw point diagnostic & 28 & 0.7857 & 0.6429 & -0.1429 [-0.3214, 0.0] & converted lower \\
Official KuJiaLe InternNav val-unseen & Success rate & 99 & 0.5253 & 0.4848 & -0.0404 [-0.1515, 0.0707] & noMDL lower \\
Official KuJiaLe InternNav val-unseen & SPL & 99 & 0.4739 & 0.4298 & -0.0441 [-0.1417, 0.0514] & noMDL lower \\
Official KuJiaLe InternNav val-unseen & Navigation error & 99 & 3.6798 & 3.6306 & -0.0492 [-0.6464, 0.5386] & noMDL lower \\
Official KuJiaLe InternNav val-unseen & Trajectory length & 99 & 6.9754 & 7.0598 & 0.0844 [-0.7366, 0.9304] & noMDL higher \\
\bottomrule
\end{tabular}%
}
\end{table*}


The paired bootstrap intervals in Table~\ref{tab:reviewer_closure_paired_ci}
quantify that comparison. On this 30-pair stress set, Gemma4 shows a small
point-hit gain, while Qwen's raw pixel diagnostic drops on converted renders.
The intervals describe frozen evidence pools.

The VLM-grounding result comes from the recorded answer and point tables.
Figures show the visible shift and provide context for those measurements.
Overlays are only guides until the exact renders have clean logs and traces.

\subsection{Material Effects and NVIDIA Baseline Boundaries}
\label{sec:material_effect_baseline}

The VLM stress result leaves a mechanism question: which cue changed the
evidence? We compare \mbox{ConvertAsset} with the NVIDIA baseline only in
checked bins.

\begin{figure*}[b]
\centering
\includegraphics[width=0.94\textwidth]{figures/fig_material_effect_baseline_qualitative.png}
\caption{Selected material-effect qualitative comparison for the four
GRScenes-covered bins. Columns show original MDL, ConvertAsset noMDL, and the
NVIDIA PreviewSurface baseline after source-MDL cells passed the clean
rerender/provenance check. Table~\ref{tab:material_effect_risk_matrix} defines
the material-effect scope.}
\label{fig:material_effect_baseline_qualitative_acl}
\end{figure*}

% Auto-generated by build_material_effect_risk_matrix.py
\begin{table*}[t]
\centering
\caption{Material-effect claim-boundary matrix for the ConvertAsset versus NVIDIA baseline. The table separates effect bins with bounded selected qualitative evidence from supplemental failure or limitation cases; it is not a population-level error-rate estimate.}
\label{tab:material_effect_risk_matrix}
\scriptsize
\resizebox{\textwidth}{!}{%
\begin{tabular}{llll}
\toprule
Effect & Evidence status & Relative interpretation & Claim allowed \\
\midrule
clearcoat & supplemental visual/PXR fail & selected NVIDIA failure; ConvertAsset target retained & selected NVIDIA failure case \\
opacity / transparency & selected qualitative ready & bounded three-condition evidence & bounded static + selected qualitative \\
emission & selected qualitative ready & bounded three-condition evidence & bounded static + selected qualitative \\
procedural texture & supplemental visual/PXR fail & both conversions lose checker evidence & diagnostic limitation case \\
normal / bump & selected qualitative ready & bounded three-condition evidence & bounded static + selected qualitative \\
displacement / height & selected qualitative ready & bounded three-condition evidence & bounded static + selected qualitative \\
\bottomrule
\end{tabular}%
}
\end{table*}


Table~\ref{tab:material_effect_risk_matrix} defines the material-effect scope.
The audit provides static evidence only for the four covered bins. Clearcoat is
one NVIDIA target-loss case, and procedural texture still has a diagnostic limit
for both converters. The figure is a visual index for those selected cases
rather than an error-rate sample.

These rows are diagnostic only. They mark cues rather than rates or full MDL
semantics.

\subsection{Stack-Entry Evidence and Official Scene Stability}
\label{sec:embodied_sanity}

The final check asks whether converted scenes can still enter an executable
embodied stack. The official
KuJiaLe \texttt{val\_unseen} route covers 99 paired InternNav/DualVLN episodes
across three local scenes. Means are SR 0.5253/0.4848, SPL 0.4739/0.4298, NE
3.6798/3.6306, and TL 6.9754/7.0598.
The paired confidence intervals in Table~\ref{tab:reviewer_closure_paired_ci}
cross zero for SR, SPL, NE, and TL. This yields a limited result: the selected
official stack can load, render, and evaluate converted scenes end-to-end.
Rollout stills and videos are qualitative checks tied to the 99-episode paired
run.

\begin{table*}[t]
\centering
\footnotesize
\setlength{\tabcolsep}{2pt}
\begin{tabular}{@{}>{\raggedright\arraybackslash}p{0.30\textwidth}>{\raggedright\arraybackslash}p{0.30\textwidth}>{\raggedright\arraybackslash}p{0.30\textwidth}@{}}
\toprule
Metric & Original MDL & noMDL \\
\midrule
Scenes & 3 & 3 \\
Successful runs & 9; 0 fail & 9; 0 fail \\
Ready time (s) & 13.95 [11.64, 16.28] & 14.12 [12.31, 16.24] \\
FPS & 2.36 [2.01, 2.70] & 2.42 [2.08, 2.73] \\
\bottomrule
\end{tabular}
\caption{Official KuJiaLe / InteriorAgent scene load/render stability. Aggregate original/noMDL rows only; intervals are bootstrap 95\% CIs. The NVIDIA official-scene row is omitted until matching converted scenes pass the same load/render check.}
\label{tab:official_scene_performance_summary}
\end{table*}


Table~\ref{tab:official_scene_performance_summary} reports load/render stability
for the same three official scenes. The 18 required paired runs all
finish in fresh launches for both asset families (18/18). Warm-up ready time is
similar for the two families. This check covers load/run stability only; speed
is outside the reported measurement. We omit the NVIDIA row until its converted scenes pass this load/render
check.
